Este trabajo describe el desarrollo de un sistema web de autogestión para el SENA, diseñado para administrar la asignación de instructores y el seguimiento de aprendices en etapa productiva.

Inicialmente, el proyecto carecía de una adecuada organización, lo que dificultaba su mantenimiento. Sin embargo, al incorporar patrones de diseño como Repository y arquitecturas en capas, se logró una significativa mejora estructural.

Los resultados incluyen un código más comprensible, reducción del 40\% en tiempos de respuesta en operaciones clave y facilidad para futuras ampliaciones. Este estudio demuestra que la aplicación de buenas prácticas transforma un proyecto caótico en una solución robusta y sostenible.